% =========================================================================
% defense.tex — 学位论文开题答辩骨架（基于 beamerthemeDefense 通用主题，脱敏骨架）
% -------------------------------------------------------------------------
% 25 页全帧型骨架：帧型结构、组件几何与实战 deck 完全一致，内容为占位文本，
% 占位即填写说明。实战原版（真实题目/姓名/图片）在
% Mac ~/Documents/code/sc_resilience_geo/docs/开题报告/beamer-dmu/defense.tex（实战原版，不入仓库）
% 编译：latexmk -xelatex defense.tex（xelatex 跑两遍定位 tikz overlay）
% =========================================================================
\documentclass[aspectratio=169,11pt]{beamer}
\usepackage{booktabs}
\usepackage{tikz}
\usetikzlibrary{arrows.meta,positioning,matrix}
\usetheme{Defense}

% 底部进度条条目名（与 \section 顺序一致；当前节深蓝高亮、其余灰色）
\DefenseSetProgress{选题背景与意义,文献综述与研究不足,研究内容与创新点,研究难点,研究方案与工作计划}
% 转场页副述（英文逗号分隔，与 \section 顺序一致；不配置则无副述行）
\DefenseSetSectionDesc{第一节一句话副述占位,第二节一句话副述占位,第三节一句话副述占位,第四节一句话副述占位,第五节一句话副述占位}

\title{论文题目（在此填写，过长可断行\\第二行占位）}
\subtitle{学位论文开题报告}
\author{答辩人：××× \and 导师：××× 教授}
\institute{大连海事大学 · 交通运输工程学院}
\date{2026 年 × 月 × 日 · 中国大连}

\AtBeginSection[]{\frame{\sectionpage}}

\begin{document}

%---------------------------------------------------------------
% 封面
%---------------------------------------------------------------
\begin{coverframe}
  \titlepage
\end{coverframe}

%---------------------------------------------------------------
% 目录
%---------------------------------------------------------------
\begin{frame}{目录}
  \begin{outlinelist}
    \outlineitem{选题背景与意义}{一句话副述：研究背景与量化目标（占位）。}
    \outlineitem{文献综述与研究不足}{一句话副述：几条脉络与缺口（占位）。}
    \outlineitem{研究内容与创新点}{一句话副述：几部分内容、几个关键问题（占位）。}
    \outlineitem{研究难点}{一句话副述：难点与对策（占位）。}
    \outlineitem{研究方案与工作计划}{一句话副述：方法体系、技术路线与进度（占位）。}
  \end{outlinelist}
\end{frame}

%===============================================================
\section{选题背景与意义}
%===============================================================

% 帧型：三卡片并列（columns[0.325] × 3 + card）
\begin{frame}{选题背景：（占位）主标题写成一个论断}
  \small
  \begin{columns}[T]
    \begin{column}{0.325\textwidth}
      \begin{card}{背景维度一}
        \footnotesize
        背景维度一的一句话概括（占位）：
        \begin{itemize}
          \item 关键事实或数据，数字用 \textcolor{accent}{强调色} 标出
          \item 要点占位文本
          \item 要点占位文本
        \end{itemize}
      \end{card}
    \end{column}
    \begin{column}{0.325\textwidth}
      \begin{card}{背景维度二}
        \footnotesize
        背景维度二的一句话概括（占位）：
        \begin{itemize}
          \item 要点占位文本
          \item 要点占位文本
          \item 要点占位文本
        \end{itemize}
      \end{card}
    \end{column}
    \begin{column}{0.325\textwidth}
      \begin{card}{背景维度三}
        \footnotesize
        背景维度三的一句话概括（占位）：
        \begin{itemize}
          \item 要点占位文本
          \item 要点占位文本
          \item 要点占位文本
        \end{itemize}
      \end{card}
    \end{column}
  \end{columns}
\end{frame}

% 帧型：三卡片（研究目标类内容：enumerate + itemize 混排）
\begin{frame}{研究目的与应用价值}
  \footnotesize
  \begin{columns}[T]
    \begin{column}{0.325\textwidth}
      \begin{card}{研究目标}
        \begin{enumerate}
          \item 目标一（占位）：构建/攻克/产出的动作句式
          \item 目标二（占位）：一句话点名技术难点
          \item 目标三（占位）：一句话点名产出形态
        \end{enumerate}
      \end{card}
    \end{column}
    \begin{column}{0.325\textwidth}
      \begin{card}{量化验证目标}
        在（占位）情景下：
        \begin{itemize}
          \item 指标一改善 \textcolor{accent}{20\%以上}
          \item 指标二代价控制在 \textcolor{accent}{5\%以内}
        \end{itemize}
        \vskip0.6mm
        一句话点出价值落点（占位）
      \end{card}
    \end{column}
    \begin{column}{0.325\textwidth}
      \begin{card}{应用价值}
        \begin{itemize}
          \item \textcolor{accent}{价值一}：一句话占位
          \item \textcolor{accent}{价值二}：一句话占位
          \item \textcolor{accent}{价值三}：一句话占位
        \end{itemize}
      \end{card}
    \end{column}
  \end{columns}
\end{frame}

%===============================================================
\section{文献综述与研究不足}
%===============================================================

% 帧型：综述页（point 分组要点 + warn 不足警示）
\begin{frame}{文献综述Ⅰ：（占位）第一条脉络名}
  \footnotesize
  \begin{point}{三类现行方法（分组标题占位）}
    \begin{itemize}
      \item \textcolor{primary}{方法类一}：代表文献（年份）——一句话概括贡献与局限
      \item \textcolor{primary}{方法类二}：代表文献（年份）；代表文献（年份）
      \item \textcolor{primary}{方法类三}：代表文献（年份）——一句话概括
    \end{itemize}
  \end{point}
  \begin{warn}{不足}
    ① 不足一占位\quad ② 不足二占位\quad ③ 不足三占位
  \end{warn}
\end{frame}

\begin{frame}{文献综述Ⅱ：（占位）第二条脉络名}
  \footnotesize
  \begin{point}{几条研究脉络（占位）}
    \begin{itemize}
      \item \textcolor{primary}{脉络一}：代表文献（年份）——一句话概括
      \item \textcolor{primary}{脉络二}：代表文献（年份）；代表文献（年份）
      \item \textcolor{primary}{脉络三}：代表文献（年份）——一句话概括
      \item \textcolor{primary}{脉络四}：模型谱系一句话概括（占位）
      \item \textcolor{primary}{脉络五}：核心结论一句话（占位）
    \end{itemize}
  \end{point}
  \begin{warn}{不足}
    该组脉络的共性不足，一句话点名（占位）
  \end{warn}
\end{frame}

\begin{frame}{文献综述Ⅲ：（占位）第三条脉络名}
  \footnotesize
  \begin{point}{三个层面（占位）}
    \begin{itemize}
      \item \textcolor{primary}{层面一}：代表文献（年份）——一句话概括
      \item \textcolor{primary}{层面二}：代表文献（年份）；代表文献（年份）
      \item \textcolor{primary}{层面三}：代表文献（年份）——一句话概括
    \end{itemize}
  \end{point}
  \begin{warn}{不足}
    该组脉络的共性不足，一句话点名（占位）
  \end{warn}
\end{frame}

% 帧型：研究缺口 + 切入点（point 编号项 \item[Ⅰ] + band 横幅）
\begin{frame}{研究缺口与本课题切入点}
  \small
  \vspace*{-2.5mm}
  \begin{point}{几点研究不足（呼应前几条脉络）}
    \begin{itemize}
      \item[Ⅰ] 不足一占位：\textcolor{accent}{强调缺口关键词}
      \item[Ⅱ] 不足二占位：\textcolor{accent}{强调缺口关键词}
      \item[Ⅲ] 不足三占位：\textcolor{accent}{强调缺口关键词}
    \end{itemize}
  \end{point}
  \vskip-1.2mm
  \begin{band}
    {\footnotesize\color{accent}本课题切入点}\quad 方法一（占位）$\;\to\;$ 方法二（占位）$\;\to\;$ 方法三（占位）$\;\to\;$ 验证手段（占位）
  \end{band}
\end{frame}

%===============================================================
\section{研究内容与创新点}
%===============================================================

% 帧型：研究框架总览（tikz stage 节点链 + 虚线反馈弧）
\begin{frame}{研究框架总览：（占位）一句话概括闭环}
  \begin{center}
    \begin{tikzpicture}[
        stage/.style={fill=primary!5, rounded corners=1.5pt,
                      text width=22mm, minimum height=17mm, align=center, inner sep=1.5mm},
        arr/.style={-{Stealth[length=2.6mm]}, line width=1.5pt, accent}]
      \node[stage] (s1)
        {{\color{accent}\bfseries\large ①}\\[0.3mm]
         \textcolor{primary}{\bfseries 模块一}\\[0.5mm]
         {\scriptsize\color{muted} 方法关键词（占位）}};
      \node[stage, right=2.5mm of s1] (s2)
        {{\color{accent}\bfseries\large ②}\\[0.3mm]
         \textcolor{primary}{\bfseries 模块二}\\[0.5mm]
         {\scriptsize\color{muted} 方法关键词（占位）}};
      \node[stage, right=2.5mm of s2] (s3)
        {{\color{accent}\bfseries\large ③}\\[0.3mm]
         \textcolor{primary}{\bfseries 模块三}\\[0.5mm]
         {\scriptsize\color{muted} 方法关键词（占位）}};
      \node[stage, right=2.5mm of s3] (s4)
        {{\color{accent}\bfseries\large ④}\\[0.3mm]
         \textcolor{primary}{\bfseries 模块四}\\[0.5mm]
         {\scriptsize\color{muted} 方法·算法（占位）}};
      \node[stage, right=2.5mm of s4] (s5)
        {{\color{accent}\bfseries\large ⑤}\\[0.3mm]
         \textcolor{primary}{\bfseries 模块五}\\[0.5mm]
         {\scriptsize\color{muted} 产出形态（占位）}};
      \draw[arr] (s1) -- (s2);
      \draw[arr] (s2) -- (s3);
      \draw[arr] (s3) -- (s4);
      \draw[arr] (s4) -- (s5);
      \draw[arr, dashed] (s5.south) to[bend left=20]
        node[below, font=\scriptsize, text=muted]{反馈闭环说明（占位）} (s1.south);
    \end{tikzpicture}
  \end{center}
\end{frame}

% 帧型：图文页（左 point×2 + 右图；论文图放 figures/ 后按下行引用）
\begin{frame}{研究内容Ⅰ：（占位）内容标题}
  \footnotesize
  \begin{columns}[T]
    \begin{column}{0.44\textwidth}
      \begin{point}{子内容一标题（占位）}
        \begin{itemize}
          \item \textcolor{primary}{数据}：数据库名与范围（占位）
          \item \textcolor{primary}{方法与目标}：方法一句话，目标一句话（占位）
        \end{itemize}
      \end{point}
      \begin{point}{子内容二标题（占位）}
        \begin{itemize}
          \item \textcolor{primary}{变量}：变量与来源一句话（占位）
        \end{itemize}
      \end{point}
    \end{column}
    \begin{column}{0.52\textwidth}
      \begin{center}
        % 论文图放 figures/ 目录后取消注释（矢量 PDF 优先）：
        % \includegraphics[width=\linewidth]{figures/figure1}
        \fbox{\parbox[c][50mm][c]{0.92\linewidth}{\centering\color{muted} 论文图占位\\[0.5mm]
          \scriptsize \texttt{figures/figure1}}}

        \vskip0.8mm
        {\scriptsize\color{muted} 图注一句话（占位）}
      \end{center}
    \end{column}
  \end{columns}
\end{frame}

% 帧型：图文页（右图用 height 约束的变体）
\begin{frame}{研究内容Ⅱ：（占位）内容标题}
  \footnotesize
  \begin{columns}[T]
    \begin{column}{0.50\textwidth}
      \begin{point}{子内容一标题（占位）}
        \begin{itemize}
          \item \textcolor{primary}{主体}：对象与角色一句话（占位）
          \item \textcolor{primary}{规则与目标}：规则一句话，目标一句话（占位）
        \end{itemize}
      \end{point}
      \begin{point}{子内容二标题（占位）}
        \begin{itemize}
          \item \textcolor{primary}{维度框架}：维度名缩写一句话（占位）
          \item \textcolor{primary}{聚合口径}：双口径/尾部风险一句话（占位）
        \end{itemize}
      \end{point}
    \end{column}
    \begin{column}{0.46\textwidth}
      \begin{center}
        % \includegraphics[height=44mm]{figures/figure2}
        \fbox{\parbox[c][44mm][c]{0.92\linewidth}{\centering\color{muted} 论文图占位\\[0.5mm]
          \scriptsize \texttt{figures/figure2}}}

        \vskip0.8mm
        {\scriptsize\color{muted} 图注一句话（占位）}
      \end{center}
    \end{column}
  \end{columns}
\end{frame}

% 帧型：纯文字方法页（point 内含公式与灰色符号图例）
\begin{frame}{研究内容Ⅲ：（占位）模型/方法标题}
  \small
  \begin{point}{子内容一标题（占位）}
    \begin{itemize}
      \item \textcolor{primary}{目标}：三维/多维协同一句话（占位）
      \item \textcolor{primary}{约束}：关键约束一句话（占位）
    \end{itemize}
    \vskip0.6mm
    $\min F=(f_1,f_2)$,\ \ $\min_x f_1(x)+\mathbb{E}_\omega[Q(x,\omega)]$\par
    {\scriptsize\color{muted} $f_1$ 目标一占位 $f_2$ 目标二占位；$x$ 一阶段决策；$Q$ 二阶段成本（符号图例占位）}
  \end{point}
  \begin{point}{子内容二标题（占位）}
    \begin{itemize}
      \item \textcolor{primary}{算法}：全局探索 + 场景降维（占位）一句话
      \item 输出形态一句话（占位）
    \end{itemize}
  \end{point}
\end{frame}

\begin{frame}{几个关键问题}
  \small
  \begin{point}{关键问题 1：（占位）问题名}
    一句话点出从什么到什么（占位）
  \end{point}
  \begin{point}{关键问题 2：（占位）问题名}
    一句话点出框架与口径（占位）
    \[ \mathcal{R}=\tfrac{1}{T\,Q_0}\textstyle\int_0^{T}\! Q(t)\,\mathrm{d}t,\qquad \mathbb{E}[\mathcal{R}],\ \min_\omega \mathcal{R}(\omega) \]
    {\scriptsize\color{muted} 公式排版示范：$Q(t)$ 性能曲线；$T$ 时长；$Q_0$ 基准（符号图例占位）}
  \end{point}
  \begin{point}{关键问题 3：（占位）问题名}
    一句话点出模型与算法（占位）
  \end{point}
\end{frame}

% 帧型：三卡片（创新点）
\begin{frame}{主要创新点}
  \footnotesize
  \begin{columns}[T]
    \begin{column}{0.325\textwidth}
      \begin{card}{创新点 1}
        \textcolor{primary}{创新点名（占位）}

        \vskip1mm
        与现有方法的差异一句话（占位）；弥补什么缺陷一句话（占位）
      \end{card}
    \end{column}
    \begin{column}{0.325\textwidth}
      \begin{card}{创新点 2}
        \textcolor{primary}{创新点名（占位）}

        \vskip1mm
        突破什么、纳入什么，形成什么可计算形态——一句话占位
      \end{card}
    \end{column}
    \begin{column}{0.325\textwidth}
      \begin{card}{创新点 3}
        \textcolor{primary}{创新点名（占位）}

        \vskip1mm
        模型 + 算法 + 输出形态一句话占位
      \end{card}
    \end{column}
  \end{columns}
\end{frame}

%===============================================================
\section{研究难点}
%===============================================================

% 帧型：难点对策表（booktabs 两列表 + \addlinespace）
\begin{frame}{预期难点与对策}
  \footnotesize
  \begin{tabular}{@{}p{0.30\textwidth}@{\quad}p{0.66\textwidth}@{}}
    \toprule
    \textcolor{primary}{\textbf{难点}} & \textcolor{primary}{\textbf{对策}} \\
    \midrule
    难点一占位 & 对策一句话占位；备选手段一句话占位 \\
    \addlinespace[0.7mm]
    难点二占位 & 对策一句话占位；校验方式一句话占位 \\
    \addlinespace[0.7mm]
    难点三占位 & 对策一句话占位；降偏差手段一句话占位 \\
    \addlinespace[0.7mm]
    难点四占位（如 NP-hard） & 算法框架占位；下界评估与加速手段占位 \\
    \addlinespace[1.2mm]
    难点五占位（如交叉验证） & 团队分工占位；先公开数据后实证的两步走占位 \\
    \bottomrule
  \end{tabular}
\end{frame}

%===============================================================
\section{研究方案与工作计划}
%===============================================================

% 帧型：双卡片（方法模块 + 数据源）
\begin{frame}{研究方法体系与数据基础}
  \footnotesize
  \begin{columns}[T]
    \begin{column}{0.475\textwidth}
      \begin{card}{方法模块}
        \begin{itemize}
          \item \textcolor{primary}{模块一}：方法一句话（占位）
          \item \textcolor{primary}{模块二}：方法与算法一句话（占位）
          \item \textcolor{primary}{模块三}：仿真/验证一句话（占位）
        \end{itemize}
      \end{card}
    \end{column}
    \begin{column}{0.475\textwidth}
      \begin{card}{数据源}
        \begin{itemize}
          \item 公开数据库一（占位）
          \item 公开数据库二（占位）
          \item 事件/文本数据（占位）
          \item 企业合作数据（占位，注明脱敏协议）
        \end{itemize}
      \end{card}
    \end{column}
  \end{columns}
\end{frame}

% 帧型：技术路线（tikz cnode+bnode 上下贴合并排五列）
\begin{frame}{技术路线：五阶段递进}
  \begin{center}
    \begin{tikzpicture}[
        cnode/.style={fill=primary, text=white, font=\footnotesize\bfseries,
                       text width=23mm, minimum height=6.5mm, align=center,
                       inner xsep=1mm, inner ysep=0mm},
        bnode/.style={draw=primary!28, fill=primary!5, line width=0.5pt,
                        font=\scriptsize, text width=23mm, minimum height=21mm,
                        align=center, inner sep=1.2mm},
        arr/.style={-{Stealth[length=2.6mm,width=2mm]}, line width=1.1pt, accent}]
      \node[cnode] (c1) {①\; 阶段一};
      \node[cnode, right=3.3mm of c1] (c2) {②\; 阶段二};
      \node[cnode, right=3.3mm of c2] (c3) {③\; 阶段三};
      \node[cnode, right=3.3mm of c3] (c4) {④\; 阶段四};
      \node[cnode, right=3.3mm of c4] (c5) {⑤\; 阶段五};
      \node[bnode, below=0.5pt of c1] (b1) {工作要点\\[0.3mm] 工作要点\\[0.3mm] 工作要点};
      \node[bnode, below=0.5pt of c2] (b2) {工作要点\\[0.3mm] 工作要点\\[0.3mm] 工作要点};
      \node[bnode, below=0.5pt of c3] (b3) {工作要点\\[0.3mm] 工作要点\\[0.3mm] 工作要点};
      \node[bnode, below=0.5pt of c4] (b4) {工作要点\\[0.3mm] 工作要点\\[0.3mm] 工作要点};
      \node[bnode, below=0.5pt of c5] (b5) {工作要点\\[0.3mm] 工作要点\\[0.3mm] 工作要点};
      \draw[arr] (c1.east) -- (c2.west);
      \draw[arr] (c2.east) -- (c3.west);
      \draw[arr] (c3.east) -- (c4.west);
      \draw[arr] (c4.east) -- (c5.west);
    \end{tikzpicture}

    \vspace{2.5mm}
    {\scriptsize\color{muted}五阶段逐级递进：每阶段的输出即下一阶段的输入（一句话概括闭环，占位）}
  \end{center}
\end{frame}

% 帧型：研究基础（双图占位 + point）
\begin{frame}{研究基础}
  \footnotesize
  \begin{center}
    % 前期工作图放 figures/ 后取消注释：
    % \includegraphics[height=34mm]{figures/figure3}\,
    % \includegraphics[height=34mm]{figures/figure4}
    \fbox{\parbox[c][34mm][c]{0.36\linewidth}{\centering\color{muted} 前期图一占位}}\,
    \fbox{\parbox[c][34mm][c]{0.36\linewidth}{\centering\color{muted} 前期图二占位}}

    \vskip0.6mm
    {\scriptsize\color{muted} 左：图一说明（占位）；右：图二说明（占位）}
  \end{center}
  \vskip1mm
  \begin{point}{探索性工作}
    前期工作逐条列出（占位）：已完成的分析/原型/数据面板；团队支撑；本人实践经验——每条一句话
  \end{point}
\end{frame}

% 帧型：进度表（booktabs 三列表）
\begin{frame}{学位论文工作计划}
  \footnotesize
  \begin{tabular}{@{}l@{\quad}p{0.46\textwidth}@{\quad}p{0.22\textwidth}@{}}
    \toprule
    \textcolor{primary}{\textbf{起讫时间}} & \textcolor{primary}{\textbf{主要研究内容}} & \textcolor{primary}{\textbf{预期成果}} \\
    \midrule
    2026.09--2026.12 & 阶段一内容一句话占位 & 学术论文 1 篇 \\
    \addlinespace[1mm]
    2027.01--2027.04 & 阶段二内容一句话占位 & 学术论文 1 篇 \\
    \addlinespace[1mm]
    2027.05--2027.10 & 阶段三内容一句话占位 & 学术论文 1 篇 \\
    \addlinespace[1mm]
    2027.11--2028.02 & 阶段四内容一句话占位 & 学术论文 1 篇 \\
    \addlinespace[1mm]
    2028.03--2028.06 & 学位论文撰写、预答辩与修改完善 & 论文外审稿 \\
    \bottomrule
  \end{tabular}

  \vskip1mm
  \begin{center}\footnotesize\color{muted}
  经费保障：（占位：项目名称 + 配套说明）
  \end{center}
\end{frame}

%---------------------------------------------------------------
% 结束页
%---------------------------------------------------------------
\begin{coverframe}
  \vfill
  \begin{center}
    {\usebeamercolor[fg]{title}\bfseries\fontsize{26}{34}\selectfont 恳请各位专家批评指正}
    \vskip6mm
    \begin{tikzpicture}
      \fill[primary] (0,0) rectangle (11mm,1pt);
      \fill[accent] (12mm,0) rectangle (22mm,1pt);
    \end{tikzpicture}
    \vskip8mm
    {\usebeamercolor[fg]{institute}\small 答辩人：××× \quad 导师：××× 教授 \quad 2026 年 × 月}
  \end{center}
  \vfill
\end{coverframe}

\end{document}
